\documentclass[border=10pt]{standalone}
\usepackage[utf8]{inputenc}
\usepackage{nacal}
\usepackage{pgfplots}
\usepackage{tikz-3dplot}
\usepackage{tikz}
\usetikzlibrary{matrix,decorations.pathreplacing}
\usetikzlibrary{calc,angles,positioning,intersections,quotes,decorations.markings,arrows.meta}
\usepackage{tkz-euclide}

\begin{document}
  \def\Rojo{red!50!black}
  \def\RojoOscuro{red!20!black}
  \def\RojoClaro{red!90!black}
  \def\Verde{green!50!black}
  \def\VerdeOscuro{green!30!black}
  \def\VerdeClaro{green!50!gray}
  \def\Azul{blue!50!black}
  \def\AzulOscuro{blue!35!black}
  \def\AzulClaro{blue!60!gray}
  \def\Datos{green!40!gray}
  \def\Ajuste{blue!50!red}
  \def\Residuo{black!25!gray}
  \def\Regresor{blue!75!white}

  % n=5 (n-k=3): t_c = 3,18  =>  theta_c = 28,56 grados
  \def\ThetaC{28.56}
  \def\Rotulo{$n=5$ \quad ($n-k=3$), \quad $t_c=3{,}18$\\[2pt]
    abertura $\theta_c=28{,}6^\circ$\\[3pt]
    $|t|=1{,}73$: \textbf{no se rechaza} $H_0$}

  \begin{tikzpicture}
    [scale=1.15,
     axis/.style={-,gray,very thin},
     vector/.style={-{Stealth[length=2.6mm, width=1.8mm]},very thick},
     guide/.style={dashed,thin,gray}]
  
    \path[use as bounding box] (-2.85,-3.75) rectangle (2.85,2.75);
  
    \coordinate (O) at (0,0);
  
    % región crítica: cono doble de semiabertura \ThetaC
    \fill[\Rojo!12] (O) -- ({-\ThetaC}:2) arc[start angle={-\ThetaC}, end angle={\ThetaC}, radius=2] -- cycle;
    \fill[\Rojo!12] (O) -- ({180-\ThetaC}:2) arc[start angle={180-\ThetaC}, end angle={180+\ThetaC}, radius=2] -- cycle;
    \draw[\Rojo!60,very thin] (O) -- ({\ThetaC}:2);
    \draw[\Rojo!60,very thin] (O) -- ({-\ThetaC}:2);
    \draw[\Rojo!60,very thin] (O) -- ({180-\ThetaC}:2);
    \draw[\Rojo!60,very thin] (O) -- ({180+\ThetaC}:2);
  
    \draw[gray!50,very thin] (O) circle (2);
  
    % los dos ejes del plano de desviaciones
    \draw[axis] (-2.35,0) -- (2.35,0);
    \node[anchor=north west,gray,yshift=-0.5mm] at (-2.35,0)
      {$\mathcal{L}(\boldsymbol{x}-\boldsymbol{\mathop{\overline{x}}})$};
    \draw[guide] (0,-2.3) -- (0,2.3);
  
    % abertura del cono (dibujada hacia abajo, donde no hay nada más)
    \draw[\Rojo,thin] (0.95,0) arc[start angle=0, end angle={-\ThetaC}, radius=0.95];
    \node[\Rojo,font=\small,inner sep=1pt] at ({-\ThetaC/2}:1.2) {$\theta_c$};
  
    % el vector de datos en desviaciones, siempre a 45 grados
    \draw[vector,\Datos] (O) -- (45:2)
      node[anchor=south west,xshift=-0.5mm]
      {$\boldsymbol{y}-\boldsymbol{\mathop{\overline{y}}}$};
    \draw[\Verde,thin] (1.55,0) arc[start angle=0, end angle=45, radius=1.55];
    \node[\Verde,font=\small,anchor=south west,inner sep=1pt] at (22.5:1.6) {$45^\circ$};
  
    \node[anchor=north,align=center,font=\small] at (0,-2.45) {\Rotulo};
  \end{tikzpicture}
\end{document}
